\section{الفصل~\ref{sec:recognition}. التعرّف}

سرعة التعرّف، وبنية المراحل الأولى، والقراءة الخطية للجواب، والتعلم من استمرارية الزمن مقيسة لدى الإنسان والقرد؛ واختزال السلسلة كلها إلى عمليتين والتعلم من الفيديو مختبَران على نماذج الكتاب؛ أما تفسيرات الخدع فتتراوح بين الراسخ الأسس والمختلف فيه.

{\footnotesize
\setlength{\tabcolsep}{3pt}\renewcommand{\arraystretch}{1.15}
\begin{longtable}{>{\raggedright}p{0.5cm}>{\raggedright}p{4.0cm}>{\raggedright}p{5.6cm}>{\raggedright}p{2.3cm}>{\raggedright\arraybackslash}p{1.7cm}}
\toprule
م & القضية & التجربة وما قيس & المصدر & الحالة \\
\midrule\endfirsthead
\toprule
م & القضية & التجربة وما قيس & المصدر & الحالة \\
\midrule\endhead
1 & المقارنة بالنموذج بكسلًا ببكسل مع الإزاحة لا تصيب بأفضل من رمي عملة؛ وزوج المراحل~\eqref{eq:scstage} يتعرف على الشكل في أي مكان & \foreignlanguage{english}{\texttt{models/\allowbreak recognition\_models.py}}، ”شبكية“ من $24$ نقطة، $4000$ محاولة: النموذج $0.49$، $0.51$، $0.52$ عند نسبة نقاط مقلوبة $0$، $0.1$، $0.2$؛ وزوج $S$ و$C$: $1.00$، $0.86$، $0.70$ & استنتاج في الفصل & نموذج \\
2 & يُتعرَّف على حيوان في صورة في $\sim150\ms$، بمرور واحد عبر نحو عشر مراحل من $10\text{--}20\ms$ لكل منها & الإنسان، مهمة ”هل يوجد حيوان أم لا“ على صور معروضة $20\ms$: الكمونات المستثارة للجوابين تفترق بعد $\sim150\ms$. القرد: زمن الاستجابة الأولى يزداد من مرحلة إلى أخرى (\foreignlanguage{english}{V1} أولًا، ثم \foreignlanguage{english}{V2} و\foreignlanguage{english}{V4}). عدد المراحل و”صفر نبضات أو نبضة واحدة لكل مرحلة“ استنتاج من هذه الأرقام & \cite{Thorpe1996,Schmolesky1998} & مُقاس \\
3 & المرحلة $S$ ”و“ على الأجزاء، والمرحلة $C$ الأكبر على المواضع (القضية~\ref{prop:conveyor}) & القط، القشرة البصرية الأولى: الخلايا البسيطة تستجيب لحافة باتجاه معيّن في مكان معيّن، والمعقدة للاتجاه نفسه في منطقة أوسع. تسجيل داخل خلوي للخلايا المعقدة: الاستجابة لشريطين قريبة لدى خلايا كثيرة من أكبر الاستجابتين، وفي المتوسط أقرب ما تكون إلى عملية \foreignlanguage{english}{max}. الأكبر عبر التثبيط المتبادل هو القضية~\ref{prop:wta}، والقسمة على النشاط الكلي مراجعة & \cite{HubelWiesel1962,Lampl2004,Carandini2012} & مُقاس \\
4 & كلما صعدنا في السلسلة كبرت التراكيب وتعقدت، وقلّ اعتماد الجواب على الموضع & القرد، تبسيط المنبهات إلى ”السمة الحرجة“: خلايا \foreignlanguage{english}{V4} والقشرة الصدغية السفلية الخلفية تستجيب لتراكيب معقدة من الشكل واللون والنسيج في مجال صغير؛ وفي القشرة الصدغية السفلية الأمامية المجال أكبر. وتناوب $S$ و$C$ في النموذج يعيد إنتاج هذه الصورة & \cite{KobatakeTanaka1994,Riesenhuber1999} & مُقاس \\
5 & الشبكات الالتفافية تكرر هذا التناوب وتتنبأ باستجابات المراحل العليا أفضل من غيرها؛ والشيء يُقرأ خطيًا من ترددات مجموعة من العصبونات & شبكة مدرَّبة على التعرّف بلا مواءمة مع الدماغ: الطبقة العليا تتنبأ باستجابات عصبونات القشرة الصدغية السفلية لدى القرد، والوسطى باستجابات \foreignlanguage{english}{V4}. مصنِّف خطي على نشاط $\sim100$ خلية مختارة عشوائيًا في نوافذ تبدأ من $12.5\ms$ يتعرف على الشيء وفئته مع اختلاف الموضع والحجم & \cite{Fukushima1980,Yamins2014,Hung2005} & مُقاس \\
6 & قاعدة هيب مع أثر~\eqref{eq:traceinvariance} تعلّم اللاتغيّر من فيديو بلا أسماء إذا لم يكن الأثر أقصر من $\sim20\ms$ & فكرة السمات البطيئة والقاعدة مع أثر نظرية. \foreignlanguage{english}{\texttt{models/\allowbreak recognition\_models.py}}، عصبونا $C$، $16$ مدخلًا، $10$ تشغيلات: التوقف بين الأشياء $0.3\,\text{s}$. عند $\tau_{\text{tr}}=5\ms$ التطابق مع الشكل $0.52$ في المتوسط، وعند $10\ms$ $0.56$؛ وعند $20$ و$50$ و$200\ms$ $1.00$ في كل التشغيلات (عند $30\ms$ يخطئ تشغيل واحد من عشرة) & \cite{Foldiak1991,WiskottSejnowski2002} & نموذج \\
7 & في الدماغ يُتعلَّم اللاتغيّر من استمرارية الزمن & القرد: استُبدل الشيء خفيةً أثناء قفزة العين إلى مكان معيّن من المجال؛ فتغيّر لاتغيّرُ الموضع لدى عصبونات القشرة الصدغية السفلية وفق الاستبدال، تغيرًا دالًّا بعد ساعة. أما أنها القاعدة الرئيسية لتعلم البصر ففرضية & \cite{LiDiCarlo2008} & مُقاس \\
8 & الحركة: كواشف الاتجاه، ثم زوج ”و“/”أو“ نفسه على مناطق كبيرة & كواشف الاتجاه في شبكية الأرنب؛ وفي المنطقة \foreignlanguage{english}{MT} لدى القرد كل العصبونات الـ$110$ المسجلة انتقائية للاتجاه، والاستجابة للحركة أقوى منها في \foreignlanguage{english}{V1} & \cite{Barlow1965,Albright1984} & مُقاس \\
9 & المراحل العليا ترسل التنبؤ إلى أسفل، والخطأ يصعد إلى أعلى & النموذج يفسر التأثيرات خارج المجال الكلاسيكي في القشرة البصرية الأولى؛ وبعض المشاهدات يتفق معه (مراجعة)، ولم يُثبت بوصفه البنية العامة لخط الأنابيب & \cite{RaoBallard1999,Keller2018} & فرضية \\
10 & الصور الصعبة تتطلب وصلات راجعة وعدة دورات & القرد: في الصور ”الصعبة“ يتكوّن القرار في القشرة الصدغية السفلية متأخرًا بنحو $\sim30\ms$؛ وهذه الاستجابات المتأخرة تتنبأ بها الشبكة الأمامية تنبؤًا ضعيفًا، والشبكات ذات الوصلات الراجعة أو الأعمق تنبؤًا أفضل & \cite{Kar2019} & مُقاس \\
11 & تزييف غير ملحوظ في البكسلات يخدع الشبكة؛ والبصر البشري أمتن بكثير، لكنه ليس محصنًا & على الشبكات: اضطراب صغير مختار خصيصًا يغيّر الجواب؛ ومثال ”الباندا $\leftarrow$ الجِبّون“ من العمل الثاني. ولدى الإنسان عند العرض القصير تزيح التزييفاتُ التي تنتقل جيدًا بين الشبكات الاختيارَ & \cite{Szegedy2014,Goodfellow2015,Elsayed2018} & مُقاس \\
12 & خدع التوقع: الدخل غير الموثوق يتنحى للتوقع، فالباهت يتحرك ”أبطأ“، والفستان يُرى على نحو مختلف (القسم~\ref{sec:illusions}) & المحزّزات الأقل تباينًا تبدو أبطأ حركة. استطلاع لـ$1401$ شخص: يُرى الفستان في ثلاثة أشكال ثابتة، وتلميح الإضاءة يبدّل اللون؛ ولدى $\sim13\,000$ مشاهد يحسم الأمرَ افتراضُ الإضاءة. والنموذج البايزي مع توقع الحركات البطيئة يفسر عددًا من خدع الحركة & \cite{Thompson1982,LaferSousa2015,Wallisch2017,Weiss2002,Purves2011} & فرضية \\
13 & ضبط الكسب: الشلال، والميل، والتباين؛ وشبكة هيرمان ليست تثبيط الجيران & كواشف الاتجاه في شبكية الأرنب تخفض ترددها عند الحركة الطويلة، وبعد التوقف تهبط دون الخلفية. خدعة الميل يعيد إنتاجها نموذج بالقسمة على نشاط الجيران. وتعديلات شبكة هيرمان التي لا تغيّر استجابة عصبونات مخرج الشبكية تزيل البقع، فرُفض التفسير بتثبيط الجيران & \cite{BarlowHill1963,Schwartz2009,SchillerCarvey2005} & مُقاس \\
14 & تعويض التأخيرات: الشيء المتحرك يبدو متقدمًا على الومضة & الخدعة مقيسة. والتجارب النفسية الفيزيائية تناقض كلًّا من مدّ الحركة إلى الأمام وفرق التأخيرات؛ واقتُرح تفسير بأثر رجعي، بحسب الأحداث خلال $\sim80\ms$ بعد الومضة. والخلاف لم يُحسم & \cite{Nijhawan1994,EaglemanSejnowski2000} & فرضية \\
15 & ”الثعابين“ الساكنة تدور: فرق التأخيرات~\eqref{eq:latency} تقرؤه كواشف الاتجاه & الخدعة تعمل بلا لون أيضًا؛ وأزواج العناصر المتجاورة تولّدها منفردة؛ وعصبونات قشرة القرد الانتقائية للاتجاه تعطي استجابة موجَّهة لهذه الأزواج الساكنة. ولدى الإنسان يطلق الدورانَ القفزاتُ الدقيقة للعين والرمش & \cite{Conway2005,OteroMillan2012} & مُقاس \\
16 & الصور المزدوجة: تثبيط متبادل، وتعب، وضجيج؛ والتبدلات يسببها الضجيج في الغالب & نموذج مجموعات عصبونات متنافسة: التبدلات فيه يسببها \emph{الضجيج}، ولا تحدث بدونه؛ ويفسر ازدياد تكرار التبدلات مع شدة المنبه. والتعب هنا دوره أصغر & \cite{MorenoBote2007} & فرضية \\
17 & بعض الخدع نتيجة للمسألة التي تتعلمها الآلة & الشبكة المدرَّبة على التنبؤ بالإطار التالي من الفيديو ”ترى“ دوران الحلقات الساكنة ولا تراه في الصور الضابطة؛ والشبكات المخصصة لتنقية الصور وزيادة حدتها تنخدع بخدع اللون والسطوع كالإنسان & \cite{Watanabe2018,GomezVilla2020} & مُقاس \\
\bottomrule
\end{longtable}}
